\section{System Design}
\label{sec:design}

\system selects which Agents' histories deserve protection (O1),
which KV regions within those histories should be retained (O2), and
how much HBM replication may occupy as foreground demand changes (O3). 
Figure~\ref{fig:system-overview} shows our four cooperating components.
\emph{Tracker} estimates remaining workflow work and recovery urgency.
\emph{Router} chooses execution workers and replica destinations.
\emph{KV Master} selects Agent histories and KV regions, evaluates
candidate placements, and manages replica retention and reclamation.
\emph{Adapter} adjusts the HBM replication budget to foreground demand. Rank-local
Connectors execute transfers and report completion to KV Master.

\begin{figure}[t]
 \centering
 \includegraphics[width=\columnwidth]{figures/system-overview.pdf}
 \caption{\system: Four-component architecture. Each workflow node is one
 invocation. Letters identify Agents, and parentheses give invocation
 ordinals, B and C are each invoked twice.
 Each host contains one TP=2 DP worker. Each rank's HBM holds execution
 and replica regions. Rank-matched Connectors copy KV between hosts;
 each host's DRAM connects to its own Connectors. Blue paths carry
 KV; dashed paths carry control.}
 \label{fig:system-overview}
\end{figure}


\subsection{Tracker}
\label{sec:tracker}

% 中文：一个 Agent 通过多步完成其任务，每步对应一次模型调用。Tracker 按调用粒度
% 记录执行和历史使用关系，按 Agent 剩余工作的粒度估计其对工作流完成时间的影响。
An Agent completes its assigned task through a sequence of steps, each
corresponding to one model invocation. Tracker records execution and
history use at invocation granularity, and estimates remaining Agent
work to determine workflow urgency.

% 中文：本文考虑 Agent 任务 DAG 预先已知且固定的工作流。Agent 集合及依赖关系已知，
% 但每个 Agent 内部的调用次数和完成时间未知。
We target workflows with a fixed Agent-task DAG known before execution.
Agent identities and dependencies are fixed, while each Agent's
invocation count and completion time are unknown.

\subsubsection{Remaining-Time Prediction}
\label{sec:tracker-prediction}

% 中文：以 (w,a,r) 标识工作流 w 中 Agent a 的第 r 次调用，当依赖满足后，workflow runtime 将这次 invocation 提交给 serving 系统调度执行。
% 后续调用可以复用兼容前缀的已选择 KV；调用结束不等于历史失效，直到 Agent 结束。
Tracker identifies invocation $r$ of Agent $a$ in workflow $w$ as
$v=(w,a,r)$. Once its dependencies are satisfied, the workflow runtime
submits the invocation to the serving system for scheduling. Successive
invocations of the same Agent may reuse selected KV for a compatible
history prefix. Completing an invocation does not end the Agent's
history; Tracker tracks that history until the Agent finishes its task.

% 中文：每次调用提交时，Tracker 根据本次调用的信息预测 Agent 距离完成任务还需多久。
% 一次调用只预测一次，下一次调用到来时再更新，直到 Agent 结束。
At each invocation submission, Tracker predicts how long the Agent will
take to finish its task using the invocation's input and the Agent's
observed history. Each invocation triggers one prediction, which is
updated at each subsequent invocation until the Agent finishes its task.

% 中文：将预测建模为基于上下文的回归问题。文本编码器提取任务描述、本次新增输入
% 和近期已接受输出的语义表示，再与提交时已知的模型、角色、输入长度、已完成调用
% 次数、历史调用的长度和耗时，以及 Agent 已运行时间结合。
% 回归器直接预测 Agent 的剩余完成时间，供后续工作流 Slack 计算使用。
We formulate this prediction as context-conditioned regression. A text
encoder represents the Agent's task description, newly appended input,
and recent accepted outputs. A regressor combines this semantic
representation with features available at submission: the model and
Agent role, input length, completed-call count, previous calls' lengths
and timings, and elapsed Agent time. It directly predicts Agent $a$'s
remaining completion time, $\widehat R_{a,r}$, for workflow slack
computation.

% 中文：训练样本将每次调用提交时的特征与该次提交到 Agent 任务结束的实际时间配对。
Training examples pair these submission-time features with the observed
duration from submission of invocation $r$ to completion of Agent $a$'s
task.

% 中文：尚未启动的 Agent 使用同模型、同角色历史任务的时长先验 D。
% 首次调用提交时，以该先验初始化剩余时间估计 R，并记录提交时间 tau。
% 每次调用触发的预测返回后，以新预测替换 R，tau 始终取该预测对应的调用提交时间。
For an unstarted Agent, Tracker uses a duration prior $\bar D_a$ from
historical tasks with the same model and role. Its first invocation
submission initializes $\widehat R_a$ to this prior and $\tau_a$ to the
submission time. Each prediction $\widehat R_{a,r}$ replaces
$\widehat R_a$ when available, with $\tau_a$ set to invocation $r$'s
submission time.

\subsubsection{Slack Computation}
\label{sec:tracker-slack}

% 中文：注册工作流时，runtime 从固定编排图导出所有 Agent 任务的工作流 ID、Agent ID
% 和前驱 Agent ID 列表；边 a→b 表示 b 等待 a 完成任务后才能开始。
% runtime 报告任务完成事件，Tracker 移除已完成节点，得到剩余 DAG。
% 该 DAG 每个节点是一个尚未完成的 Agent 任务，其内部可以有多次 invocation。
At workflow registration, the runtime exports each Agent task's workflow
ID, Agent ID, and predecessor Agent IDs from the fixed orchestration graph.
An edge $a\to b$ means that $b$ waits for Agent $a$ to finish its task.
The runtime reports task completion, allowing Tracker to remove finished
tasks and maintain the residual DAG $G_t=(A_t,E_t)$. Each node in this
graph is one unfinished Agent task, which may contain multiple invocations.

% 中文：算法以剩余 Agent DAG、时长先验、已启动 Agent 的最新估计及提交时间、误差裕量为输入。
% 未启动 Agent 保留完整时长先验；已启动 Agent 的估计换算到当前时刻。
% 正向计算最早完成时间，反向计算最晚完成时间。
% 后向遍历以预测的工作流完成时间为终点；每个 Agent 的剩余工作只计算一次。
% pred 和 succ 分别表示剩余前驱与后继；g 是经校准的 Slack 高估误差裕量。
Algorithm~\ref{alg:slack-estimation} computes guarded Agent slack using
forward and backward critical-path passes~\cite{kelley1959cpm}.
Unstarted Agents retain their full duration priors, while estimates for
started Agents are adjusted to the current time $t$. The algorithm counts
each Agent's remaining work once and uses the predicted workflow finish
$T$ as the backward-pass endpoint. The two passes compute earliest and
latest finishes, $EF_a$ and $LF_a$. Here $\operatorname{pred}(a)$ and
$\operatorname{succ}(a)$ denote residual predecessors and successors,
and $g_a$ is a calibrated upper quantile of slack overestimation.

% 中文：所需时长估计缺失、过期或无效时，整个工作流暂不授予正 Slack。
% 正向遍历：EF = 本节点剩余时间 + 最晚完成的前驱；后向遍历：
% LF = 后继最晚开始时间的最小值，汇点从 T 开始回推。
% 最终 Slack 为 max(0, LF - EF - g)；两次遍历的复杂度均为 O(|A_t|+|E_t|)。
\par\medskip\noindent\begin{minipage}{\columnwidth}
\raggedright
\refstepcounter{designalgorithm}\label{alg:slack-estimation}
\textbf{Algorithm \thedesignalgorithm: Compute Agent slack}
\par\smallskip
\begin{tabularx}{\columnwidth}{@{}r@{\quad}X@{}}
\toprule
 & \textbf{Input:} $G_t=(A_t,E_t)$; Agent states; priors $\bar D_a$;
 estimates $(\widehat R_a,\tau_a)$; guards $g_a$; time $t$ \\
 & \textbf{Output:} guarded slack $\boldsymbol{s}=\{s_a:a\in A_t\}$ \\
\midrule
1 & $s_a\gets0,\quad \forall a\in A_t$ \\
2 & \textbf{for} $a\in A_t$ \textbf{do} \\
3 & \quad \textbf{if} $a$ is unstarted \textbf{then} $d_a\gets\bar D_a$ \\
4 & \quad \textbf{else} $d_a\gets\widehat R_a-(t-\tau_a)$ \\
5 & \textbf{if} any required estimate is missing or invalid, or
 any $d_a\le0$, \textbf{then return} $\boldsymbol{s}$ \\
6 & $\pi\gets\textsc{TopologicalOrder}(G_t)$ \\
7 & \textbf{for} $a$ in $\pi$ \textbf{do} \\
8 & \quad $EF_a\gets d_a+
 \max\bigl(\{0\}\cup\{EF_p:p\in\operatorname{pred}(a)\}\bigr)$ \\
9 & $T\gets\max\bigl(\{0\}\cup\{EF_a:a\in A_t\}\bigr)$ \\
10 & \textbf{for} $a$ in reverse $\pi$ \textbf{do} \\
11 & \quad $LF_a\gets
 \min\bigl(\{T\}\cup\{LF_b-d_b:b\in\operatorname{succ}(a)\}\bigr)$ \\
12 & \quad $s_a\gets\max\{0,LF_a-EF_a-g_a\}$ \\
13 & \textbf{return} $\boldsymbol{s}$ \\
\bottomrule
\end{tabularx}
\end{minipage}\par\medskip

% 中文：调用继承所属 Agent 的 Slack。
Invocation $v=(w,a,r)$ inherits its Agent's slack, $s_v=s_a$.

% 中文：预测和 Slack 更新在 Controller 中异步执行，调用调度和 token 生成无需等待。
Prediction and slack updates run asynchronously in the controller, so
invocation dispatch and token generation do not wait for them.

\subsection{Router and Connector}
\label{sec:router}

% 中文：Router 根据队列负载和前缀 KV 的位置调度调用，与此同时 Router 还会决定如果 Failure 发生时，该 Agent 调用 retry 到哪个DP Replica
% 这一选择也决定了如果要备份该调用对应的 KV，备份在哪里。
% retry 目标位于执行节点的受保护故障域之外，选择时考虑负载、可用内存和复制路径竞争。
Router schedules invocations based on queue load and prefix locality.
It also selects a retry DP replica for each Agent invocation in case of
failure. If the invocation's KV is selected for backup, it is placed at
this retry target. The target lies outside the execution worker's
protected failure domain and is chosen based on load, available memory,
and copy-path contention.

% 中文：在该目标上，主机内存（Peer DRAM）保存受保护 Agent 的完整 KV 检查点；GPU 内存
% （Peer HBM）额外保存 KV Master 从其前 L 层中选出的 KV 块。容量不足时，Router 将缺口
% 返回 KV Master。
At this target, host memory (\emph{Peer DRAM}) holds a protected Agent's
complete KV checkpoint, and GPU memory (\emph{Peer HBM}) additionally
holds the KV blocks that KV Master selects from its first $L$ layers.
If a placement does not fit, Router returns the capacity shortfall to
KV Master.

% 中文：各 TP rank 的 Connector 执行 KV 复制并向 KV Master 报告完成状态。
% 故障后 Router 使旧执行失效，并将调用重定向到 retry 目标；Engine 直接复用目标
% HBM 中的 KV，Connector 将 DRAM 中的 KV 搬入 GPU。目标成为执行节点后，
% Router 为后续备份重新选择 Peer。
Rank-local Connectors execute KV copies and report completion to KV
Master. After a failure, Router fences the old execution and redirects
the invocation to the retry target. The engine reuses resident Peer HBM
KV, while Connectors move DRAM-resident KV into GPU memory.
If that target becomes the execution worker, Router selects a new peer
for subsequent backups.

% 中文：上述路由与传输机制处理已选状态；接下来由 KV Master 决定保护哪些
% Agent 历史，以及保留其中哪些 KV 区域。
These routing and transfer mechanisms operate on selected state.
KV Master determines which Agent histories and which KV regions to
protect, as described next.

\subsection{KV Master}
\label{sec:master}

% 中文：KV Master 决定每个 retry 目标保存什么。在 Agent 之间，它选择哪些历史在 Peer DRAM
% 中保存完整的 KV 检查点（O1）；在受保护的 Agent 内部，它选择哪些 KV 同时放入 Peer HBM，
% 使恢复不必等待 DRAM 传输（O2）。检查点以固定的 256 token 粒度推进。
KV Master decides what each retry target holds. Across Agents, it selects
the histories that receive a complete KV checkpoint in Peer DRAM (O1).
Within a protected Agent, it selects the KV that Peer HBM also holds, so
that recovery does not wait for DRAM transfers (O2). Checkpoints advance
in fixed chunks of 256 tokens.

\subsubsection{Agent Selection}
\label{sec:agent-protection}

% 中文：令 T_a 为 Agent a 被中断的调用通过完整重建重试所需的延迟。将 T_a 和 Tracker 的
% Slack s_a 代入式 (1)，得到 Slack 无法掩盖的重试延迟 E_a。
Let $T_a$ be the delay of retrying Agent $a$'s interrupted invocation by
full reconstruction. Applying Equation~\ref{eq:isolated_penalty} with
$R_i=T_a$ and Tracker's guarded slack $\sigma_i=s_a$ gives the retry delay
that slack cannot hide,
\begin{equation}
 E_a=\bigl[T_a-s_a\bigr]^+.
 \label{eq:exposed-delay}
\end{equation}
% 中文：当 E_a > λ 时，Master 为 a 在 Peer DRAM 中保存完整的 KV 检查点。λ 是可以容忍的
% 工作流延迟：未掩盖延迟不超过 λ 的 Agent 失败后直接重建；λ = 0 时，重建超过 Slack 的
% Agent 都受保护。DRAM 放不下全部入选 Agent 时，Master 按 E_a 从大到小接纳。T_a 随历史
% 增长，s_a 随工作流推进变化，因此 Master 在调用边界和 Slack 更新时重新判断。
Master keeps a complete KV checkpoint of $a$ in Peer DRAM when
$E_a>\lambda$. The threshold $\lambda$ is the workflow delay that the
operator tolerates. Agents whose unhidden delay does not exceed $\lambda$
are reconstructed after a failure, and $\lambda=0$ protects every Agent
whose reconstruction exceeds its slack. If Peer DRAM cannot hold every
selected Agent, Master admits them in decreasing $E_a$. Because $T_a$
grows with the history and $s_a$ changes as the workflow advances, Master
repeats the test at invocation boundaries and slack updates.

% 中文：图 X(a) 将该规则用于图 2(b) 中的 Agent。A1 的重建落在 Slack 内；取 λ = 17 s 时，
% A2 和 A4 获得 DRAM 检查点，而 A3 的 16.3 s 不超过 λ，失败后重建。
Figure~\ref{fig:master-selection}(a) applies this rule to the Agents of
Figure~\ref{fig:agent_value}(b). A1's reconstruction fits within its
slack. With $\lambda=17$\,s, A2 and A4 receive DRAM checkpoints, whereas
A3's 16.3\,s does not exceed $\lambda$, so A3 is reconstructed if it fails.

% 中文（图示）：构造示例。(a) 使用图 2(b) 的未掩盖延迟，λ 为示意取值；(b) 块布局为示意。
\begin{figure}[t]
 \centering
 \includegraphics[width=\columnwidth]{figures/master-selection.pdf}
 \caption{Selection by KV Master. (a) Retry delays
 not hidden by slack for the Agents of Figure~\ref{fig:agent_value}(b);
 with $\lambda=17$\,s, A2 and A4 receive complete DRAM checkpoints.
 (b) One protected Agent with $K=2$ and $L=2$ (dashed line). Every
 checkpointed block is in DRAM (D), the two most accessed blocks of L1 and
 L2 are also in HBM (H), and the tail of at most 256 tokens after the
 checkpoint is rebuilt during recovery.}
 \label{fig:master-selection}
 \Description{Top: horizontal bars of unhidden retry delay for Agents A1
 to A4, with 0, 32.6, 16.3 and 18.0 seconds, and a dashed threshold line
 at 17 seconds; A2 and A4 exceed it. Bottom: a grid of layers L1 to L4 and
 blocks B0 to B5. B0 and B3 of L1 and B1 and B3 of L2 are marked H; all
 other blocks are marked D. A hatched tail column is marked rebuilt, and a
 dashed line below L2 marks depth L.}
\end{figure}

\subsubsection{Checkpoint Granularity}
\label{sec:master-checkpoint}

% 中文：受保护 Agent 的检查点按固定的 256 token 推进：每新接受 256 个 token，Connector
% 将它们所有层的 KV 复制到 Peer DRAM，Master 随后发布新的检查点位置。故障后，retry 目标
% 恢复检查点，并重建检查点之后至多 256 个已接受 token 的 KV，与 checkpoint-and-replay
% 恢复一致。固定粒度把重建限制在一次 256 token 的 Prefill 内，并使每个 token 的 KV 只以
% 大块传输复制一次。
A protected Agent's checkpoint advances in fixed chunks of 256 tokens.
Once 256 new tokens are accepted, Connectors copy their KV for all layers
to Peer DRAM, and Master publishes the new checkpoint position. After a
failure, the retry target restores the checkpoint and rebuilds the KV of
the at most 256 accepted tokens after it, as in checkpoint-and-replay
recovery~\cite{strati2024dejavu,cao2026lumen}. The fixed chunk bounds the
rebuild to a single 256-token prefill, and each token's KV is copied only
once, as part of a large transfer.

\subsubsection{HBM Placement}
\label{sec:region-selection}

% 中文：从 DRAM 恢复可以与重建重叠：重建逐层进行，GPU 计算前面的层时，Connector 可以
% 传输后面层的 KV（图 3(b)）。仍有两部分传输暴露：第 1 层之前没有计算可以掩盖其传输；
% 当一层 KV 的传输慢于一层的重建时，等待逐层累积。Master 因此把受保护 Agent 前 L 层中
% 每层的 Top-K 块同时放入 Peer HBM。
Restoring from DRAM overlaps with the rebuild. The rebuild proceeds layer
by layer, so Connectors load a later layer's KV while the GPU computes
earlier layers (Figure~\ref{fig:kv_value}(b)). Two parts of the transfer
remain exposed. No computation precedes layer~1 to hide its transfer, and
when a layer's KV loads more slowly than a layer rebuilds, waiting
accumulates across layers. Master therefore also places the Top-$K$
blocks of each of the first $L$ layers of a protected Agent in Peer HBM.

\paragraph{Top-$K$ blocks.}
% 中文：O2 表明稀疏注意力的读取集中在少数块上。Master 按引擎在正常执行中记录的原生选择
% 次数对每层的块排序，保留访问最多的 K 个。K 按模型固定，由离线测得的访问集中度
% （图 3(a)）确定。重建前 L 层时直接读取 HBM 中的 Top-K 块；对后续层，Connector 按层序
% 先传 Top-K 块，其余 KV 随后传输；选中但不在 GPU 中的块按需从 DRAM 读取。
O2 shows that sparse attention concentrates its reads on a few blocks.
Master ranks each layer's blocks by the native selection counts that the
engine records during execution and keeps the $K$ most accessed. $K$ is
fixed per model, from its access concentration measured offline
(Figure~\ref{fig:kv_value}(a)). The rebuild of the first $L$ layers reads
their Top-$K$ blocks in HBM. For later layers, Connectors load the Top-$K$
blocks first in layer order and the rest of the KV afterwards. Selected
blocks that are not yet in GPU memory are fetched from DRAM on demand.

% 中文：图 X 示意这一效应。(a) 先从 DRAM 传完全部 KV 再重建，整个传输都暴露。(b) 按层序
% 先传 Top-K 块，传输与重建重叠，但 Top-K 块的传输慢于重建时每层仍要等待。(c) L1–L3 的
% Top-K 块在 HBM 中，重建立即开始，后续 Top-K 块都在需要之前到达，不再等待；其余 KV 在
% 调用就绪后继续传输。
Figure~\ref{fig:recovery-pipeline} illustrates the effect. Loading the
whole KV from DRAM before the rebuild exposes the entire transfer (a).
Loading Top-$K$ blocks first in layer order overlaps the transfer with the
rebuild, but every layer still waits when its Top-$K$ blocks load more
slowly than a layer rebuilds (b). With the Top-$K$ blocks of L1--L3 in
HBM, the rebuild starts at once, later Top-$K$ blocks arrive before they
are needed, and no layer waits (c). The rest of the KV continues to load
after the invocation is ready.

% 中文（图示）：示意图，不对应特定模型或实测：8 层，时间以一层重建时间为单位；Top-K 块占
% 每层 KV 的一半，传输需 1.4 个单位，即带宽低于 B*。
\begin{figure}[t]
 \centering
 \includegraphics[width=\columnwidth]{figures/recovery-pipeline.pdf}
 \caption{Illustrative recovery of one protected Agent with eight layers. Each
 dark segment rebuilds one layer for the tail of at most 256 tokens. Top-$K$
 blocks are half of each layer's KV and load more slowly than a layer
 rebuilds. (a) The whole KV loads from DRAM before the rebuild. (b) Top-$K$
 blocks load first in layer order. (c) The Top-$K$ blocks of L1--L3 are in
 HBM. The rest of the KV loads after the Top-$K$ blocks.}
 \label{fig:recovery-pipeline}
 \Description{Three two-lane timelines, H2D above GPU. (a) The GPU waits
 while all KV loads, then rebuilds eight layers. (b) Top-K blocks load
 layer by layer and every layer waits briefly before it rebuilds; the rest
 of the KV loads afterwards. (c) With the first three layers' Top-K blocks
 in HBM, the eight layers rebuild without waiting while loads continue.}
\end{figure}

\paragraph{Depth $L$.}
% 中文：L 取决于单个恢复实际获得的有效 H2D 带宽。传输很少能打满链路，并发恢复还会共享
% 链路，因此该带宽低于链路峰值，需要在预期的并发恢复下离线测量。令 B* 为一层 Top-K 块
% 的传输时间恰好等于一层重建时间时的带宽。带宽不低于 B* 时只需保存第 1 层（L = 1）；
% 低于 B* 时传输落后于重建，L 随带宽下降迅速增大（图 Y）。部署据此为每个模型设定 L。
$L$ depends on the effective host-to-device bandwidth that one recovery
obtains. Transfers rarely saturate the link, and concurrent recoveries
share it, so this bandwidth lies below the link's peak and is measured
offline under the expected concurrent recoveries. Let $B^*$ be the
bandwidth at which a layer's Top-$K$ blocks load in the time a layer
rebuilds. At or above $B^*$, holding layer~1 suffices, so $L=1$
(Figure~\ref{fig:depth-bandwidth}). Below $B^*$, loads fall behind the
rebuild, and $L$ grows quickly as the bandwidth drops. The deployment sets
$L$ for each model from this profile.

% 中文（图示）：示意图，32 层，与图 X 采用相同的规则。
\begin{figure}[t]
 \centering
 \includegraphics[width=\columnwidth]{figures/depth-bandwidth.pdf}
 \caption{Illustrative depth $L$ versus effective host-to-device bandwidth
 for a 32-layer model. Bandwidth is normalized to $B^*$, at which a layer's Top-$K$
 blocks load in one layer's rebuild time. (a) Smallest $L$ for which no
 layer waits. (b) Waiting without HBM, in layer rebuild times. Below $B^*$
 (dashed), loads fall behind the rebuild and $L$ grows as the bandwidth
 drops.}
 \label{fig:depth-bandwidth}
 \Description{Two stacked panels over bandwidth from 0 to 2 times B-star.
 (a) Depth L falls in steps from 31 near zero bandwidth to 1 at B-star and
 stays 1 above. (b) On a logarithmic scale, waiting without HBM falls from
 several hundred layer rebuild times near zero bandwidth to one at B-star,
 then declines slowly.}
\end{figure}

\paragraph{HBM admission.}
% 中文：每个受保护 Agent 占用 L·K 个块的 HBM。只有从 DRAM 恢复的延迟 T_a^D 仍超过 Slack
% 的 Agent 才能从 HBM 获益。Master 按 [T_a^D − s_a]^+ 从大到小将 Agent 接纳到 Adapter
% 给出的 HBM 配额中，跳过该值为 0 的 Agent；放不下的 Agent 保留 DRAM 检查点，Master 将其
% 报告给 Adapter，作为增大配额的信号。
Each protected Agent occupies $L\cdot K$ blocks of HBM. HBM helps only
Agents whose recovery from DRAM, taking $T^{D}_a$, still exceeds their
slack. Master admits protected Agents to Adapter's HBM quota in
decreasing $[T^{D}_a-s_a]^+$ and skips those for which it is zero. Agents
that do not fit keep their DRAM checkpoint, and Master reports them to
Adapter as its signal for growing the quota (Section~\ref{sec:capacity}).

% 中文：图 X(b) 展示一个受保护 Agent 的布局：每个已检查点的块都在 DRAM 中；K = 2、L = 2
% 时，L1 和 L2 中访问最多的两个块同时在 HBM 中；检查点之后的尾部在恢复时重建。
Figure~\ref{fig:master-selection}(b) shows the resulting layout for one
protected Agent. Every checkpointed block is in DRAM. With $K=2$ and
$L=2$, the two most accessed blocks of L1 and L2 are also in HBM, and the
tail after the checkpoint is rebuilt during recovery.

\subsubsection{Replica Management}
\label{sec:master-online}
\label{sec:execution}

% 中文：副本增量更新。每推进一个检查点，DRAM 只追加新的 256 token 块；Master 重新排序
% Top-K，只把新进入 Top-K 的块复制到 HBM，离开的块在正在进行的读取结束后释放。同一 Agent
% 的后续调用扩展同一前缀，因此副本跨调用有效；Agent 完成任务后释放其副本。
Replicas are updated incrementally. Each checkpoint appends only the new
256-token chunk to DRAM. Master then re-ranks the Top-$K$ blocks, copies
only blocks that enter the Top-$K$ set to HBM, and releases blocks that
leave it once active reads finish. Replicas remain valid across
invocations, because later invocations of the same Agent extend the same
prefix, and they are released when the Agent finishes its task.

% 中文：复制是异步的。副本目录记录每个块的层、token 区间、TP 分片和位置，所有分片到达后
% 才发布检查点或 HBM 副本。故障后，Router 将调用重定向到 retry 目标：引擎直接使用 HBM
% 中的块，Connector 按上述顺序从 DRAM 加载其余 KV，引擎重建尾部。DRAM 检查点是完整的，
% 因此 HBM 副本只影响恢复时间，不影响正确性。
Because copies are asynchronous, a replica manifest records each block's
layer, token range, TP shards, and location, and it publishes a checkpoint
or HBM copy only after all its shards arrive. After a failure, Router
redirects the interrupted invocation to its retry target. The engine uses
HBM blocks in place, Connectors load the remaining KV from DRAM in the
order above, and the engine rebuilds the tail. Because the DRAM checkpoint
is complete, HBM copies affect recovery time but never correctness.

\subsection{Adapter}
\label{sec:capacity}

% 中文：O3 表明前台服务可让出的 HBM 随负载变化。Adapter 为每个接管 rank k 维护 HBM 副本
% 配额 B_k，KV Master 只在该配额内接纳 Agent（每个 Agent 占 L·K 个块）。配额受上限 C_k
% 约束：B^pool_k 为扣除权重和固定分配后的显存，B^serve_k 为服务高水位与执行预留，
% B^reserve_k 为增长与 workspace 余量，B^cap_k 为管理员设定的上限。
O3 shows that the HBM that foreground serving can spare changes with
load. Adapter maintains an HBM replica quota $B_k$ for each takeover rank
$k$, and KV Master admits Agents, each occupying $L\cdot K$ blocks, only
within it. The quota is bounded by
\begin{equation}
 C_k=\max\{0,\min(B^{\rm cap}_k,
 B^{\rm pool}_k-B^{\rm serve}_k-B^{\rm reserve}_k)\},
 \label{eq:ceiling}
\end{equation}
where $B^{\rm pool}_k$ is the memory left after weights and fixed
allocations, $B^{\rm serve}_k$ covers the serving high watermark and
execution reservations, $B^{\rm reserve}_k$ keeps headroom for growth and
workspace, and $B^{\rm cap}_k$ is an operator limit.

% 中文：Adapter 以 AIMD 调整 B_k。分配卡顿、预留被侵占或因显存不足的准入等待出现时，
% B_k 乘以 γ < 1；连续 m 个窗口稳定且 Master 报告有 Agent 因 HBM 不足未被接纳时，B_k
% 增加 Δ，并截断到 C_k。深度 L 与 K 不随配额变化，配额只改变可放入 HBM 的 Agent 数量。
Adapter adjusts $B_k$ by additive increase and multiplicative
decrease~\cite{chiu1989aimd}. Allocation stalls, reserve violations, or
admission waits caused by memory multiply $B_k$ by $\gamma<1$. After $m$
stable windows, Adapter adds $\Delta$ if Master reports protected Agents
deferred for HBM capacity, and clips the result to $C_k$. $L$ and $K$ stay
fixed, so the quota changes only how many Agents HBM holds.

% 中文：配额缩小时，Master 停止新的 HBM 接纳，并按 [T_a^D − s_a]^+ 从小到大释放 Agent
% 的 HBM 副本，等其读取结束后回收。DRAM 检查点是完整的，因此无需迁移数据，这些 Agent
% 只是改为从 DRAM 恢复。前台准入按实际占用而非配额判断，释放完成前不会超额分配。
When the quota shrinks, Master stops new HBM admissions and releases the
HBM copies of Agents in increasing $[T^{D}_a-s_a]^+$ once their readers
drain. Because their DRAM checkpoints are complete, no data moves, and
these Agents recover from DRAM instead. Foreground admission follows
actual occupancy rather than the quota, so memory is not overcommitted
before the release completes.
